\chapter{Methoden}
\label{cha:methoden}

In diesem Kapitel wird beschrieben, mit welchen wissenschaftlichen Methoden die in Kapitel~\ref{cha:stand_der_technik} aufgeworfenen Fragen bearbeitet werden.
Zwei Methoden werden eingesetzt.
Die Literaturrecherche dient dazu, den aktuellen Stand der Forschung zu Tor, Website-Fingerprinting und Circuit Padding aufzuarbeiten und daraus die Forschungslücke abzuleiten.
Das anschliessende Laborexperiment im Shadow-Simulator dient dazu, die aus dieser Literatur abgeleiteten Hypothesen quantitativ zu überprüfen.
Die konkrete Umsetzung dieses Versuchsaufbaus wird in Kapitel~\ref{cha:realisierung} beschrieben.
Das vorliegende Kapitel beschränkt sich auf die methodische Begründung der Vorgehensweise.

\section{TODO: WF vs Timing Angriff}

In der Zwischenpräsentation wurde bemängelt es sei zu wenige praxisorientiert. hier eingehen auf die Untershciede, den Ablauf, Begründungen und wie ich hierhin gekommen bin.

Timing Attacken auch in der Stand der Technik

\section{Methodische Einordnung der Arbeit}
\label{sec:methodische_einordnung}

Die vorliegende Bachelorarbeit ist eine Mischform aus einer konstruktiven und einer empirischen Arbeit~\cite[S.~71\,f.]{book}.
Konstruktiv ist die Arbeit, weil eine eigene Werkzeugkette aus Shadow-Steuerung, Konfigurationsgenerierung, Webserver, Datenkonvertierung und Klassifikator-Anbindung neu aufgebaut wird, um den untersuchten Effekt überhaupt messen zu können.
Empirisch ist die Arbeit, weil mit dieser Werkzeugkette ein kontrolliertes Laborexperiment durchgeführt wird, dessen Ergebnisse die formulierten Hypothesen entweder bestätigen oder widerlegen.

Die Arbeit verwendet ein deduktives Vorgehen.
Aus der Literatur, insbesondere aus den Arbeiten von Jansen und Wails~\cite{jansen2023data}, Sirinam et al.~\cite{sirinam2018deep} und Kadianakis et al.~\cite{kadianakis2021tor}, werden Hypothesen über den Einfluss von Circuit Padding auf die Klassifikationsgenauigkeit von \gls{df} abgeleitet.
Diese Hypothesen werden anschliessend an Messdaten überprüft.
Die Messung selbst erfolgt quantitativ über etablierte Klassifikationsmetriken (Accuracy, Precision, Recall, F1-Wert) sowie für den Vergleichsversuch zur Traffic-Korrelation über \gls{roc}- und \gls{auc}-Kennzahlen.

\section{Literaturrecherche}
\label{sec:literaturrecherche}

Die Literaturrecherche folgt dem in Balzert beschriebenen fünfstufigen Vorgehen aus den Schritten Was suche ich, wo suche ich, wie suche ich, was verwende ich, wie beschaffe ich~\cite[S.~215\,ff.]{book}.
Die einzelnen Stufen werden im Folgenden zusammengefasst.

\subsection{Such\-strategie und Abgrenzung des Themenfelds}
\label{sec:literaturrecherche_strategie}

Ausgangspunkt der Recherche ist das in Kapitel~1 formulierte Forschungsthema rund um den Einfluss von Circuit Padding auf Website-Fingerprinting in Shadow.
Aus dem Themenfeld werden vier inhaltliche Suchrichtungen abgeleitet, die sich in Schlagwortlisten überführen lassen: das Tor-Netzwerk und sein Bedrohungsmodell, Website-Fingerprinting-Klassifikatoren, Circuit Padding als Verteidigung sowie der Shadow-Netzwerksimulator und die zugehörige Werkzeugkette tornettools.
Pro Suchrichtung werden deutsche und englische Schlagwörter, Synonyme und verwandte Begriffe gesammelt.
Beispiele sind onion routing, traffic analysis, deep fingerprinting, padding machine, network simulation und cell direction sequence.

\subsection{Genutzte Quellen, Datenbanken und Suchmaschinen}
\label{sec:literaturrecherche_quellen}

Die Recherche stützt sich primär auf wissenschaftliche Datenbanken und nicht auf allgemeine Websuchmaschinen.
Folgende Quellen werden eingesetzt:

\begin{itemize}
    \item Google Scholar als übergreifende Suchmaschine, die zusätzlich die Vorwärts-Zitationen eines Artikels über die Funktion Cited by erschliesst.
    \item swisscovery als Bibliothekskatalog der HSLU für Lehrbücher, insbesondere für die methodologische Grundlage~\cite{book}.
    \item Primärquellen wie das offizielle Tor-Spezifikationsdokument~\cite{torspec_dirspec}, der Tor-Metrics-Datendienst~\cite{tormetricsabout} und das OnionPerf-Projekt~\cite{onionperf}.
\end{itemize}

Innerhalb der Datenbanken werden die Suchanfragen mit den gängigen booleschen Operatoren (\texttt{AND}, \texttt{OR}, \texttt{NOT}) sowie mit Phrasensuche und Feldfiltern (Titel, Abstract, Erscheinungsjahr) verfeinert.
Bei zu hoher Trefferzahl wird nach Erscheinungsjahr eingeschränkt oder ein zusätzlicher fachspezifischer Begriff hinzugefügt.
Bei zu geringer Trefferzahl werden Schlagwörter durch Synonyme ersetzt oder über das Schneeballsystem erweitert.

\subsection{Schneeball- und Vorwärtssuche}
\label{sec:literaturrecherche_schneeball}

Aus den Suchergebnissen wird eine kleine Menge von Schlüsselpublikationen identifiziert, die in der Folge als Ausgangspunkt für eine systematische Schneeballsuche dienen~\cite[S.~225]{book}.
Im rückwärtsgerichteten Schneeballverfahren werden die Literaturverzeichnisse dieser Schlüsselpublikationen, allen voran das Artefaktpapier von Jansen und Wails~\cite{jansen2023data}, vollständig durchgegangen und thematisch passende Quellen weiterverfolgt.
Mehrfach zitierte Werke werden als grundlegend eingestuft, dazu zählen unter anderem das ursprüngliche Tor-Designpapier~\cite{dingledine2004tor}, die Open-World-Analyse von Wang und Goldberg~\cite{wang2016realistically} und die DF-Originalveröffentlichung~\cite{sirinam2018deep}.
Im vorwärtsgerichteten Verfahren wird über die Cited by-Funktion von Google Scholar geprüft, in welchen jüngeren Arbeiten die identifizierten Schlüsselpublikationen aufgegriffen oder weiterentwickelt wurden.
Diese Vorgehensweise ergänzt die Schneeballsuche um neuere Publikationen, die in den älteren Quellen naturgemäss nicht enthalten sein können.

\subsection{Bewertung der Quellen}
\label{sec:literaturrecherche_bewertung}

Jede Quelle wird vor der Verwendung auf Zitierfähigkeit und Zitierwürdigkeit geprüft~\cite[S.~166\,ff.]{book}.
Zitierfähig sind Quellen, die veröffentlicht, nachvollziehbar und überprüfbar sind.
Zitierwürdig sind in dieser Arbeit ausschliesslich peer-reviewte Konferenzbeiträge, Zeitschriftenartikel, technische Berichte etablierter Institutionen sowie offizielle Spezifikations- und Projektdokumente des Tor-Projekts.
Populärwissenschaftliche Quellen wie Tageszeitungen oder Blogartikel werden nicht als Beleg für inhaltliche Aussagen verwendet.
Eine eng begrenzte Ausnahme bilden Pressemitteilungen des Tor-Projekts und einzelne Hersteller-Dokumente, die als Primärquelle für Software-Verhalten oder als Beleg für tagesaktuelle Ereignisse herangezogen werden.

Internetquellen werden zusätzlich zur Zitierwürdigkeitsprüfung mit dem Zugriffsdatum und einer lokalen Kopie versehen, sofern sie keine langfristig stabile Adresse wie eine \gls{doi} besitzen.
Damit ist sichergestellt, dass auch dann, wenn die Originalressource später verändert oder entfernt wird, der zitierte Inhalt überprüfbar bleibt~\cite[S.~168\,f.]{book}.

Bei Konflikten zwischen Quellen wird, wo möglich, auf die Primärquelle zurückgegriffen.
Sekundäre Wiedergaben werden nur dann verwendet, wenn die Primärquelle nicht zugänglich ist, was in dieser Arbeit nicht vorkommt.

\subsection{Literaturverwaltung}
\label{sec:literaturrecherche_verwaltung}

Sämtliche während der Recherche ermittelten Quellen werden in einem zentralen BibTeX-Verzeichnis (\texttt{src/thesis/references.bib}) gepflegt.
Pro Quelle werden Autor, Titel, Erscheinungsjahr, Verlag oder Konferenz, \gls{doi} oder URL und das Zugriffsdatum erfasst.
Die zugehörigen Volltexte werden parallel im Verzeichnis \texttt{sources/} unter dem jeweiligen Bibkey abgelegt.
Diese strikte Eins-zu-eins-Zuordnung von Bibkey und Volltext erleichtert eine spätere stichprobenartige Überprüfung der Zitate gegen das Original.
Im Text wird nach IEEE-Zitierstil mit Seitenangabe nach deutscher Konvention zitiert (\verb|\cite[S.~X]{key}|).

\section{Experimentelles Vorgehen}
\label{sec:experimentelles_vorgehen}

Die zentrale empirische Methode dieser Arbeit ist das kontrollierte Laborexperiment.
Dieser Abschnitt beschreibt das experimentelle Design auf methodischer Ebene.
Die konkrete technische Umsetzung folgt in Kapitel~\ref{cha:realisierung}.

\subsection{Wahl des Experimenttyps}
\label{sec:experiment_typ}

Nach Balzert lassen sich empirische Arbeiten in experimentelle und nicht-experimentelle Arbeiten unterscheiden, wobei das Experiment durch die kontrollierte Variation einer oder mehrerer unabhängiger Variablen unter konstant gehaltenen Rahmenbedingungen charakterisiert ist~\cite[S.~74\,ff.]{book}.
Innerhalb der Experimente wird zwischen Feldexperimenten in der natürlichen Umgebung und Laborexperimenten in einer kontrollierten Umgebung unterschieden.
Für die Untersuchung des Einflusses von Circuit Padding wird in dieser Arbeit ein Laborexperiment gewählt, das vollständig innerhalb des Shadow-Netzwerksimulators stattfindet.

Diese Wahl ist methodisch begründet.
Ein Feldexperiment im produktiven Tor-Netzwerk wäre nicht reproduzierbar, da Lastspitzen, Routing-Änderungen und konkurrierender Verkehr die Messung verändern, ohne dass diese Effekte vom Padding-Effekt unterschieden werden könnten~\cite[S.~1\,f.]{jansen2023data}.
Zusätzlich wäre eine zentrale Variation der Padding-Konfiguration im Live-Netzwerk nicht möglich, ohne den Betrieb fremder Tor-Nutzerinnen und -Nutzer zu beeinflussen, was ethisch nicht vertretbar wäre.
Die Shadow-Simulation erfüllt dagegen alle drei zentralen Anforderungen an ein Laborexperiment: kontrollierte Bedingungen, Reproduzierbarkeit und Wiederholbarkeit~\cite[S.~5\,f.]{netsim-atc2022}.

\subsection{Forschungsfragen und Hypothesen}
\label{sec:experiment_hypothesen}

Aus der Aufgabenstellung übernehmen.

\subsection{Variablen und Messgrössen}
\label{sec:experiment_variablen}

Das Experiment unterscheidet zwischen unabhängigen und abhängigen Variablen, wie es für ein kontrolliertes Laborexperiment gefordert ist~\cite[S.~74]{book}.

\paragraph{Unabhängige Variablen}
Drei Variablen werden gezielt variiert.
Erstens der Padding-Modus mit den drei Ausprägungen OFF, ON und Reduced, die in Tabelle~\ref{tab:padding_modi} (Kapitel~\ref{cha:realisierung}) erläutert sind.
Zweitens die Anzahl der Klassen mit den Ausprägungen 20, 80 und 280 (letzteres als Open-World mit 80 überwachten Seiten).
Drittens das Welt-Szenario mit den Ausprägungen Closed-World und Open-World.

\paragraph{Abhängige Variablen}
Gemessen werden Accuracy, Precision, Recall und F1-Wert auf einer zurückgehaltenen Testmenge.
Für den Vergleichsversuch zur Traffic-Korrelation werden zusätzlich \gls{roc}-Kurven sowie die \gls{auc} ausgewiesen, ergänzt um die True-Positive-Raten an den Schwellen $\textrm{FPR}=10^{-2}$ und $\textrm{FPR}=10^{-3}$.

\paragraph{Konstant gehaltene Rahmenbedingungen}
Alle übrigen Parameter bleiben über die Experimente einer Vergleichsgruppe konstant.
Dazu zählen die Skalierung des simulierten Netzwerks (1\,\%), der Simulator-Seed, der Seed der Trainings- und Testaufteilung, die Hyperparameter des Klassifikators sowie die Auswahl der Wikipedia-URLs.
Damit lässt sich eine beobachtete Differenz in den abhängigen Variablen kausal auf die Variation der unabhängigen Variable Padding-Modus zurückführen und nicht auf Sekundäreffekte.

\subsection{Versuchsplan}
\label{sec:experiment_plan}

Der Versuchsplan besteht aus sieben Hauptläufen, die in Tabelle~\ref{tab:versuchsplan} zusammengefasst sind.
Pro Closed-World-Konfiguration werden alle drei Padding-Modi durchlaufen, sodass ein direkter Vergleich der Klassifikationsgenauigkeit innerhalb derselben Klassenanzahl möglich ist.
Das Open-World-Experiment wird einmalig mit deaktiviertem Padding gemessen, da der Open-World-Fall vor allem zur Einordnung der Closed-World-Resultate dient.

\begin{table}[H]
\centering
\caption{Versuchsplan der Hauptexperimente.}
\label{tab:versuchsplan}
\begin{tabular}{llll}
\hline
Lauf & Klassen & Padding & Szenario \\
\hline
exp-baseline-20 & 20 & OFF & Closed-World \\
exp-padding-20 & 20 & ON & Closed-World \\
exp-reduced-20 & 20 & Reduced & Closed-World \\
exp-baseline-80 & 80 & OFF & Closed-World \\
exp-padding-80 & 80 & ON & Closed-World \\
exp-reduced-80 & 80 & Reduced & Closed-World \\
exp-openworld & 80 überwacht (280 gesamt) & OFF & Open-World \\
\hline
\end{tabular}
\end{table}

Ergänzend wird für eine Teilmenge der Konfigurationen ein Traffic-Correlation-Experiment durchgeführt (Schalter \texttt{-{}-correlation}), bei dem zusätzlich an ausgewählten Exit-Relais aufgezeichnet wird.
Dieses Experiment dient als unabhängiger Vergleichswert, weil Traffic-Korrelation einer anderen Angreiferklasse entspricht und kein vortrainiertes Modell benötigt.

\subsection{Validität, Reliabilität und Objektivität}
\label{sec:experiment_guete}

Die wissenschaftliche Güte des Experiments wird an den Kriterien Validität, Reliabilität und Objektivität gemessen~\cite[S.~21\,ff.]{book}.

\paragraph{Validität}
Die interne Validität ist hoch, weil im Shadow-Simulator alle Rahmenbedingungen ausser dem Padding-Modus konstant gehalten werden und damit eine kausale Interpretation der gemessenen Differenzen zulässig ist.
Die externe Validität, also die Übertragbarkeit der Resultate auf das Live-Tor-Netzwerk, wird gestützt durch die Cross-Dataset-Untersuchung von Jansen und Wails, in der ein in Shadow trainierter Klassifikator auf Live-Daten eine Genauigkeit von rund 86\,\% erreichte~\cite[S.~7]{jansen2023data}.
Die Übertragbarkeit ist damit nicht vollständig, aber empirisch begründet.

\paragraph{Reliabilität}
Die Reliabilität ergibt sich aus der deterministischen Ausführung von Shadow bei festem Seed.
Wiederholte Läufe mit identischer Konfiguration erzeugen byte-identische Pcaps und damit identische Trainings- und Testergebnisse.
Diese Reproduzierbarkeit wird durch den vollständigen Versionsstand aller Eingaben (Tor-Konsensdaten eines festgehaltenen Monats, ZIM-Datei, URL-Liste, Skripte) abgesichert.

\paragraph{Objektivität}
Die Auswertung erfolgt automatisiert durch das Skript \texttt{report.py}, das die Klassifikationsmetriken aus den von WFlib erzeugten Ergebnis-JSONs in eine zentrale Tabelle (\texttt{ExperimentLogs.csv}) überführt.
Eine subjektive Interpretation der Rohdaten findet nicht statt.
Die in Kapitel~\ref{cha:diskussion} dargestellten Diagramme werden ebenfalls automatisiert aus dieser Tabelle erzeugt, sodass der Weg von der Rohmessung bis zur grafischen Darstellung jederzeit reproduzierbar ist.

\subsection{Grenzen der gewählten Methode}
\label{sec:experiment_grenzen}

Die Methode hat zwei strukturelle Grenzen, die hier ehrlich benannt werden.
Erstens reduziert die Skalierung des Netzwerks auf 1\,\% die Pfaddiversität gegenüber dem Live-Netzwerk, was die Verkehrsmuster einer Klasse homogener und damit leichter klassifizierbar macht~\cite[S.~11]{jansen2023data}.
Zweitens wird der Tor Browser durch ein gepatchtes \texttt{wget2} ersetzt, das zwar denselben Satz an Ressourcen abruft, jedoch JavaScript und Browser-spezifische Verkehrsmuster nicht reproduziert~\cite[S.~5]{jansen2023data}.
Beide Einschränkungen werden bei der Interpretation der Resultate in Kapitel~\ref{cha:diskussion} berücksichtigt.

\section{Projektplan}
\label{sec:projektplan}

Abbildung~\ref{fig:projektplan} zeigt den zeitlichen Verlauf der Arbeit von Semesterbeginn am 16.~Februar~2026 bis zur Abgabe am 15.~Juni~2026. Die Meilensteine (rot) markieren fixe Termine, die Balken die fachlichen Arbeitsphasen.

\begin{figure}[H]
\centering
\resizebox{\textwidth}{!}{%
\begin{ganttchart}[
    hgrid, vgrid={*{6}{draw=none}, dotted},
    x unit=1.3mm, y unit chart=6mm,
    time slot format=isodate,
    calendar week text={\currentweek},
    milestone/.append style={fill=red!70},
    bar/.append style={fill=blue!40},
    group/.append style={fill=black!70},
  ]{2026-02-16}{2026-06-15}
  \gantttitlecalendar{month=name, week} \\
  \ganttgroup{Bachelorarbeit}{2026-02-16}{2026-06-15} \\
  \ganttbar{Literaturrecherche}{2026-02-16}{2026-03-20} \\
  \ganttbar{Konzept \& Stand der Technik}{2026-03-02}{2026-04-03} \\
  \ganttbar{Shadow-Simulation aufsetzen}{2026-03-16}{2026-04-16} \\
  \ganttbar{ML-Pipeline (WFlib, pcap$\to$npz)}{2026-04-06}{2026-05-08} \\
  \ganttbar{Experimente \& Auswertung}{2026-04-20}{2026-05-29} \\
  \ganttbar{Thesis schreiben}{2026-03-02}{2026-06-12} \\
  \ganttbar{Review \& Abgabe-Vorbereitung}{2026-06-01}{2026-06-15} \\
  \ganttmilestone{Semesterstart (16.02.2026)}{2026-02-16} \\
  \ganttmilestone{Kickoff (20.02.2026)}{2026-02-20} \\
  \ganttmilestone{Aufgabenstellung signiert (02.03.2026)}{2026-03-02} \\
  \ganttmilestone{Erste erfolgreiche Simulation (16.04.2026)}{2026-04-16} \\
  \ganttmilestone{Zwischenpr\"asentation (23.04.2026)}{2026-04-23} \\
  \ganttmilestone{Abgabe (15.06.2026)}{2026-06-15}
\end{ganttchart}
%
}
\caption{Projektplan der Bachelorarbeit mit Arbeitsphasen (blau) und fixen Meilensteinen (rot).}
\label{fig:projektplan}
\end{figure}
